\chapter{Introducción}
\label{ch:introduccion}

\section{Contexto y Relevancia Institucional}

La Universidad Nacional del Altiplano de Puno custodia, gestiona y transfiere permanentemente el patrimonio digital, académico, científico y administrativo de más de 20,000 usuarios activos entre estudiantes, investigadores, plana docente y personal directivo-administrativo. La progresiva consolidación de servicios digitales (como el Sistema de Matrícula en línea, la plataforma de Campus Virtual Moodle, el Sistema de Trámite Documentario SITRAD y los repositorios institucionales de investigación) ha desplazado los límites del campus físico hacia un ecosistema hiperconectado, distribuido y expuesto permanentemente a Internet.

\section{Caracterización de la Superficie Expuesta Universitaria}

Sin embargo, esta apertura operativa ha expandido correlativamente la superficie de exposición institucional. Las auditorías pasivas y de reconocimiento externo realizadas sobre el dominio \path{unap.edu.pe} constatan una arquitectura heterogénea donde coexisten más de un centenar de subdominios registrados, servidores web con capas de proxy inverso desactualizadas (como \texttt{nginx/1.16.1} en el entorno de matrícula), mecanismos de transporte seguro sin directivas estrictas de mitigación frente a interceptaciones (ausencia de \path{Strict-Transport-Security} --- HSTS) y una vulnerable postura de confianza en la mensajería institucional evidenciada por la inexistencia de directivas DMARC (\texttt{NXDOMAIN}) y mecanismos permisivos de Sender Policy Framework (SPF: \texttt{\string~all}).

\section{Ineficacia del Perímetro Tradicional y Nuevas Amenazas}

En este escenario, el modelo clásico de ``seguridad perimetral'' ---basado en la premisa ingenua de que los recursos dentro de la red interna son confiables por defecto--- resulta completamente ineficaz. La sofisticación de las amenazas actuales, tales como campañas de Spear-Phishing dirigidas al personal administrativo en ventanas críticas de matrícula, infecciones por ransomware que paralizan servidores de gestión pública, y técnicas de movimiento lateral basadas en el abuso de identidades válidas, exige la formulación de una estrategia integral, moderna y resiliente.

\section{Enfoque Holístico de Ciberseguridad y Gobernanza}

El presente Plan de Ciberseguridad Institucional adopta un enfoque multifacético y holístico: acoge al estándar NIST Cybersecurity Framework (CSF) 2.0 como el marco sombrilla de gobernanza; integra el modelado temprano de amenazas mediante STRIDE para los flujos de software; correlaciona los comportamientos y tácticas adversarias observadas con la taxonomía MITRE ATT\&CK Enterprise; incorpora la vigilancia contextual de Threat Intelligence; diseña barreras de contención escalonadas mediante la Cyber Kill Chain; y materializa controles estrictos bajo la doctrina de Zero Trust (\textit{``Nunca confiar, siempre verificar de manera explícita y asumir la brecha''}). Con ello, la UNAP alinea su gestión tecnológica con las exigencias del Decreto de Urgencia N° 006-2020 y asegura la continuidad irrestricta de su servicio educativo superior.
